\providecommand{\toprule}{\hline\hline}
\providecommand{\midrule}{\hline}
\providecommand{\bottomrule}{\hline\hline}
\providecommand{\addlinespace}[1][]{}
{\small
\renewcommand{\arraystretch}{1.12}
\begin{tabular}{llrrr}
\toprule
Coorte & Período & Voltou 12m: Choque $g\_{r,s}$ & Voltou 12m: Portab. $\times g\_{r,s}$ & Mudou setor: Portab. $\times g\_{r,s}$ \\
\midrule
2014 & Véspera da recessão & -0,0237 (0,0159) & +0,0220 (0,0179) & -0,0117 (0,0401) \\
2015 & Recessão & +0,0107 (0,0103) & -0,0113 (0,0124) & +0,0014 (0,0391) \\
2016 & Recessão & +0,0808*** (0,0280) & -0,0996*** (0,0304) & +0,0263 (0,0702) \\
2019 & Véspera da pandemia & +0,0059 (0,0052) & -0,0025 (0,0048) & -0,0421** (0,0169) \\
2020 & Pandemia & +0,0043 (0,0033) & -0,0036 (0,0041) & -0,0388* (0,0215) \\
2021 & Pandemia & -0,0078 (0,0220) & +0,0207 (0,0283) & -0,0288 (0,0408) \\
\bottomrule
\end{tabular}
\vspace{2pt}
\par{\footnotesize \textit{Nota:} Interação entre o índice de portabilidade da ocupação de origem (fixado em 2010) e a taxa de crescimento anual do emprego formal no par (região intermediária, divisão CNAE) no ano da demissão (g\_{r,s}).}
}
